\clearpage

\section{多体协同与共享通道}\label{sec:3}

\subsection{按车体包络判断通道清空}\label{sec:3.1}

通道导轨内侧宽约 1.02~m。早期曾用车辆中心是否越过通道出口作放行依据。观察轨迹后发现，车中心出去时，后轮或车尾仍可能占在口部。最终方案对底盘、车轮和万向球采用 0.28~m 的保守包络半径。保护区域覆盖两根导轨之间及口部，取 x = 1.90–5.10~m、|y| \ensuremath{\leq} 0.55~m 的矩形。若车身中心到矩形的最短距离不超过 0.28~m，就认为车辆仍占用通道。判空覆盖了车身驶离通道的全过程。

运输车按 R2、B1、R1、B2 的顺序放行，此外的侦察和空中观察可以并行。R2 的交付期限较早，R1 需要等到释放时刻，B2 还需要得到新鲜观察。这样的安排把入口会车问题转化为通道外等待和依次通过。次序在主回合开始前确定，实际启程时刻由净空、任务释放和观察状态共同决定。固定顺序便于检查，也减少了同时运输的机会。

另用两辆电机驱动的车单独比较放行时机。两组都看到前车在约 7.89~s 完整驶出。保留原定时段时，后车等到 32.00~s 发车。按车体实际清空放行时，后车在约 7.89~s 出发，少等约 24.11~s，并于约 17.58~s 清空通道。固定时段下的后车在 60~s 窗口结束时仍在通道内调整姿态。两个回合均未记录到通道占用重叠。在这组场地和出入口条件下，及时归还通道减少了后车的放行等待。

\subsection{入口对准和出口可用性}\label{sec:3.2}

\subsubsection*{联合准入检查}

中央通道的时间窗回答“什么时候可以用”，却不能回答“现在是否进得去”。入口附近的泊位、换乘路径和对准动作可能相互交叉。若车辆只凭预约时刻驶入，即使时间上没有重叠，也可能在入口姿态或出口泊位上发生冲突。

东口设置两个泊位和一个对准点。图示表明，换乘路径会穿过对准区，入口动作不能只按通道中心线生成。Gazebo 场景里的联合准入需要同时检查四类条件。

\figitem{0.98}{fig04.pdf}{东口泊位、固定对准点和队友换乘路径。几何关系为示意。}{4}

\begingroup\small
\setlength{\tableavailable}{\dimexpr\linewidth-4\tabcolsep\relax}
\begin{longtable}{@{}>{\raggedright\arraybackslash}p{0.28571\tableavailable}>{\raggedright\arraybackslash}p{0.35714\tableavailable}>{\raggedright\arraybackslash}p{0.35714\tableavailable}@{}}
\caption{通道准入与释放检查}\label{tab:14}\\
\toprule
\textbf{检查对象} & \textbf{通过条件} & \textbf{不满足时的处理} \\
\midrule
\endfirsthead
\multicolumn{3}{@{}l}{\small\color{muted}表~\thetable\quad （续）}\\
\toprule
\textbf{检查对象} & \textbf{通过条件} & \textbf{不满足时的处理} \\
\midrule
\endhead
\midrule
\multicolumn{3}{r@{}}{\footnotesize\color{muted}续下页}\\
\endfoot
\bottomrule
\endlastfoot
名义时间窗 & 申请区间没有与既有预约冲突 & 延后、重订或选择绕行 \\
实际占用 & 通道没有未清空的实体 & 在入口外等待，不提前记录进入 \\
入口几何 & 对准点与同伴在途移动保持安全距离 & 跳过固定对准点，退回旧两点式入口 \\
出口与目的地 & 出口方向和目的地区域可以接收 & 暂缓进入或调整目的地服务顺序 \\
\end{longtable}
\endgroup


获得预约后，车辆还须等待通道及出口具备接收条件。通道状态由预约转为使用后，以实际清空时刻重新计算释放时间，剩余名义时段可供后续合法预约。对准超时或目的地持续被占用时，可取消或重订预约。

\subsubsection*{东口对准与几何回退}

在一次 Gazebo 试验中，B2 已获通道授权，却迟迟未能进入，持有通道约 101.5~s。仿真轨迹显示，车辆在入口附近约 74~s 内只移动了 0.24~m。长时间的小幅移动出现在东口，排查因此集中到进场动作。回看车辆姿态后发现，原先边斜移边转向的动作在东口难以完成。加设固定对准点后，同场景的另一次回合记录到约 23.6~s 的占用时间。这支持继续采用入口对准。

对准点随后暴露了一个几何冲突。队友换乘路径与固定对准区域相交，最近距离约 0.497~m。有几次还出现入口目标格被判为不可通行的情况。简单保留固定对准点会把一个入口问题换成队友路径冲突，完全撤销又会回到原来的斜移转向问题。最终在程序中保留两条路径：几何检查通过时用对准点，不通过就退回旧的两点式入口。

\begingroup\small
\setlength{\tableavailable}{\dimexpr\linewidth-4\tabcolsep\relax}
\begin{longtable}{@{}>{\raggedright\arraybackslash}p{0.24336\tableavailable}>{\raggedright\arraybackslash}p{0.41654\tableavailable}>{\raggedright\arraybackslash}p{0.34010\tableavailable}@{}}
\caption{入口设计案例的推理链}\label{tab:15}\\
\toprule
\textbf{阶段} & \textbf{观察或检查} & \textbf{对设计的影响} \\
\midrule
\endfirsthead
\multicolumn{3}{@{}l}{\small\color{muted}表~\thetable\quad （续）}\\
\toprule
\textbf{阶段} & \textbf{观察或检查} & \textbf{对设计的影响} \\
\midrule
\endhead
\midrule
\multicolumn{3}{r@{}}{\footnotesize\color{muted}续下页}\\
\endfoot
\bottomrule
\endlastfoot
授权后仍无法进入 & 持有通道约 101.5~s。约 74~s 仅移动 0.24~m & 重点检查入口运动，而非继续延长名义时间窗 \\
增加固定对准点 & 同场景历史对照中占用约 23.6~s & 保留对准动作作为可行入口方案 \\
发现队友路径冲突 & 对准点与换乘线段最近约 0.497~m & 固定对准点不能无条件使用 \\
按几何条件选用 & 0.65~m 检查不通过时跳过并回退 & 同时保留改善动作与安全退路 \\
\end{longtable}
\endgroup


这些数据来自该组 Gazebo 入口对照，说明对准条件为何需要单独检查。八实体 Webots 场景按本章的车体包络核对通行。

对准点能帮助车辆在驶入口部前摆正，也可能与队友的换乘路径冲突。这个场地用 0.65~m 作为参考线段的检查阈值，测得的最近距离、回退规则和适用条件见附录~\ref{sec:C}。

\subsection{报告时效和岗位确认}\label{sec:3.3}

\subsubsection*{观察时效与六态岗位}

观察报告注明来源、内容、采样时刻和有效期，并用编号识别重复消息。接收方以到达时刻减去采样时刻，判断报告是否仍在 6~s 有效期内。Gazebo 记录排队、送达、失败和待处理的状态，只有送达且有效的报告才能用于服务或放行。

岗位状态把提议、确认、前往、到位待命、执行和退出分开。每次转换都有对应的事件记录，而不是根据画面中两车的距离推断。

\figitem{1}{fig05.pdf}{六种岗位状态及 Gazebo 历史回合中的服务事件}{5}

\subsubsection*{四阶段服务握手与跨平台落地}

一次完整服务经历四个衔接阶段，完成后记一次释放，车辆随即返回本职任务：
\begin{itemize}
\item \textbf{请求与提议}\quad 请求者说明所需服务及理由，候选服务者计算自身代价并报价。
\item \textbf{确认与激活}\quad 双方对同一提议编号确认，并在满足到位、报告有效等前置条件后激活执行。
\item \textbf{释放与复归}\quad 服务结束或不再必要时记录释放事件，释放后的车辆重新参与后续任务。
\end{itemize}

Gazebo 历史回合记录了蓝方的交接过程。B3 在 115.888~s 采样，B1 和 B2 在 116.952~s 收到，报告年龄为 1.064~s。服务以这次送达为生效依据。另一次跨队服务请求在 90.628~s 被拒绝，因为当时的任务规则只允许队内服务。图~\ref{fig:5} 保留了两类事件。

Webots 的整场控制把报告规则放到 B2 的放行条件里。等待、拒收旧帧、重新采样和获得有效帧四步都有事件时间。报告刚采集时反映的是过去一个时刻的画面，其有效期应由接收方核对。补采完成后，还要检查通道空闲以及任务已释放。Gazebo 中的岗位确认解决的是分工承诺能否履行，Webots 中的新鲜帧确认解决的是 B2 本次是否可以开走。

\subsection{协同方案的生成与筛选}\label{sec:3.4}

任务分配研究通常同时考虑任务收益、执行代价和资源约束\cite{ref12}。Khamis 等综述了多机器人任务分配的主要方法\cite{ref13}，Chakraa 等进一步讨论了优化目标与求解技术\cite{ref14}。本设计不对所有实体和任务做全排列，只按任务图、当前状态和资源生成有限候选，再在候选集里比较。

Gazebo 和离散模型把一次决策拆成四个维度，交叉组合出待检查方案：

\begingroup\small
\setlength{\tableavailable}{\dimexpr\linewidth-4\tabcolsep\relax}
\begin{longtable}{@{}>{\raggedright\arraybackslash}p{0.23000\tableavailable}>{\raggedright\arraybackslash}p{0.34000\tableavailable}>{\raggedright\arraybackslash}p{0.43000\tableavailable}@{}}
\caption{协同候选方案的四个维度}\label{tab:16}\\
\toprule
\textbf{维度} & \textbf{可选方案} & \textbf{含义} \\
\midrule
\endfirsthead
\multicolumn{3}{@{}l}{\small\color{muted}表~\thetable\quad （续）}\\
\toprule
\textbf{维度} & \textbf{可选方案} & \textbf{含义} \\
\midrule
\endhead
\midrule
\multicolumn{3}{r@{}}{\footnotesize\color{muted}续下页}\\
\endfoot
\bottomrule
\endlastfoot
运输分工 & 单车运输、双车交接 & 独立完成，或由两辆车分担运输与交接 \\
观察支援 & 固定岗位、空中观察 & 使用地面固定观察点或无人机报告 \\
运输路线 & 共享通道、外侧绕行 & 按预约经过中央，或改走外侧 \\
信息方式 & 补采、自行观察、携带信息、等待 & 根据所需信息的来源与时效选择动作 \\
\end{longtable}
\endgroup


四个维度产生 32 个待检查组合，每次决策再按当前状态筛掉一部分。预测时要看队友位置与可用承诺、报告新鲜度、空中巡检承诺、观察与通信路线，以及通道时隙；缺能力、缺承诺、信息过期或路线与当前协议不合的候选都会去掉。安全等待只保留在固定观察支援加单车运输这一种情形里，这种情形单列，不与确定预测下的普通路线混为一谈。

筛选时先检查当前平台能否提供候选所需的能力，再看任务是否有足够时间。比如用两车交接，需要第二辆车空闲、交接点可用、双方到达窗口相容。选择中央路线需要有效预约，选择绕行需要核对同一条任务期限。过滤后的方案按按期交付、新增工作量和交接空档比较，准备执行的方案还要等双方确认。

\subsection{仿生分工与军蚁活桥}\label{sec:3.5}

课程在集群智能部分概括了鸟群、蚁群等仿生群体的局部交互与群体决策思路 \cite{ref1}，并在地面无人车部分指出多车协作行为本质上是对内外部刺激的响应，需要建立数学模型并借助仿真工具验证 \cite{ref2}。群体机器人研究亦将这类机制归结为局部感知、局部交互与分布式决策 \cite{ref20}，国内已有面向无人机集群的系统论述 \cite{ref21}。本设计参考了四类群体行为，将其落到具体的判断规则上，表~\ref{tab:17} 逐项列出每个参考用在了设计的哪一步。

\begingroup\small
\setlength{\tableavailable}{\dimexpr\linewidth-4\tabcolsep\relax}
\begin{longtable}{@{}>{\raggedright\arraybackslash}p{0.32636\tableavailable}>{\raggedright\arraybackslash}p{0.33338\tableavailable}>{\raggedright\arraybackslash}p{0.34026\tableavailable}@{}}
\caption{仿生现象与本设计的关系}\label{tab:17}\\
\toprule
\textbf{理论或现象} & \textbf{从中得到的设计思路} & \textbf{在本项目中的体现} \\
\midrule
\endfirsthead
\multicolumn{3}{@{}l}{\small\color{muted}表~\thetable\quad （续）}\\
\toprule
\textbf{理论或现象} & \textbf{从中得到的设计思路} & \textbf{在本项目中的体现} \\
\midrule
\endhead
\midrule
\multicolumn{3}{r@{}}{\footnotesize\color{muted}续下页}\\
\endfoot
\bottomrule
\endlastfoot
Reynolds 的分离、对齐、聚合规则~\cite{ref7} & 用局部信息维持群体间距与方向秩序 & 取局部间距的思路处理窄口，两车让行仍由通道占用判断 \\
军蚁活桥的成本收益折中~\cite{ref6} & 有个体停下为队伍服务时，它原有工作的时间也要计入 & 在离散任务算例中估算中继岗位的机会成本 \\
社会性昆虫的响应阈值与分工~\cite{ref8} & 根据任务刺激和自身阈值决定是否响应 & 图模型中作为岗位响应扩展；Webots 对抗回合中用于撤销被虚占的通道名额（第~\ref{sec:adv.2} 节） \\
少数知情个体引导群体运动~\cite{ref9} & 专门观察的个体可把少量有时间意义的信息传给行动者 & B3 提供带时刻的观察，B2 据此等待或出发 \\
\end{longtable}
\endgroup


军蚁活桥提供了最直接的启发。一辆车暂时做中继，可能使另一辆车早些交付，这时仍要问中继车原来的工作会不会被耽误。为此构造了一个离散任务算例。比较方案时先看按期任务数，再看总逾期量，其后才比较车辆活跃工作时间与任务完成时刻。这个优先顺序把“做完几件事”和“做这些事花了多少车时”分开了。

例如一张任务卡给 R1 的期限为 50~s，给提供帮助的车辆的期限为 37~s。帮助后，R1 从 42~s 提前到 34~s，提供帮助的车辆仍在 37~s 完成，但合计工作时间由 32~s 增至 34~s。当 R1 已能按期时，节省出来的时间是否值得调车，要看项目对车时的要求。当 R1 的期限压到 34~s 时，帮助能够改变按期任务数，所以两种计分规则均选择帮助。两种计分规则在 12 组条件下各运行一次，其中 2 组的岗位选择发生变化。这个结果说明岗位互助的价值取决于任务还有多少时间可用。

\figitem{1}{fig06.pdf}{任务期限与支援成本对岗位选择的影响}{6}

\subsection{支援的机会成本}\label{sec:3.6}

要把军蚁活桥的启发量化，首先需要判断主任务与支援任务哪一项更紧迫。这里用一个效用向量 U = (按期任务数，\ensuremath{-}总逾期，\ensuremath{-}名义活跃工作时间，\ensuremath{-}完成时刻和)。比较时遇到第一项不同就定先后。R1 期限 50~s、提供者返航期限 37~s 时，不支援的向量是 (2, 0, \ensuremath{-}32, \ensuremath{-}79)，支援是 (2, 0, \ensuremath{-}34, \ensuremath{-}71)，时间在此统一为秒。两者都按期，先比较活跃工作时间，因此保留岗位。去掉活跃工作这一项后，选择完成时刻和更小的支援方案。

该算例在固定图模型经过 12 组、24 次配对运行复核，只有第三项活跃工作成本被置零，其余预测、期限判定、通信和实际验收保持不变。两处选择变化发生在 R1 期限 50~s、提供者返航期限 37~s 和 40~s，另外 10 组选择一致。

\begingroup\small
\setlength{\tableavailable}{\dimexpr\linewidth-10\tabcolsep\relax}
\begin{longtable}{@{}>{\raggedright\arraybackslash}p{0.25000\tableavailable}>{\centering\arraybackslash}p{0.13000\tableavailable}>{\centering\arraybackslash}p{0.13000\tableavailable}>{\centering\arraybackslash}p{0.13000\tableavailable}>{\centering\arraybackslash}p{0.19000\tableavailable}>{\raggedright\arraybackslash}p{0.17000\tableavailable}@{}}
\caption{R1 期限 50~s、提供者返航 37~s 时的实际运行结果}\label{tab:18}\\
\toprule
\textbf{成本处理} & \textbf{R1 交付/s} & \textbf{蓝方 1 批/s} & \textbf{蓝方 2 批/s} & \textbf{提供者返航/s} & \textbf{报价结果} \\
\midrule
\endfirsthead
\multicolumn{6}{@{}l}{\small\color{muted}表~\thetable\quad （续）}\\
\toprule
\textbf{成本处理} & \textbf{R1 交付/s} & \textbf{蓝方 1 批/s} & \textbf{蓝方 2 批/s} & \textbf{提供者返航/s} & \textbf{报价结果} \\
\midrule
\endhead
\midrule
\multicolumn{6}{r@{}}{\footnotesize\color{muted}续下页}\\
\endfoot
\bottomrule
\endlastfoot
保留活跃工作成本 & 42 & 53 & 61 & 37 & 拒绝 \\
第三项置零 & 34 & 43 & 53 & 37 & 接受并完成 \\
\end{longtable}
\endgroup


保留成本时，不提供服务使 R1 较晚交付，蓝方后续两批也相应较晚。置零后成交，R1 与蓝方交付都提前，而提供者返航仍为 37~s。该结果把仿生启发、效用排序和实际选择连了起来，但仍是固定确定性图模型算例。该结果用于解释岗位规则。

当 R1 期限收紧到 34~s 时，两种成本处理都接受服务，实际 R1 交付 34~s、蓝方两批 43~s 和 53~s、提供者返航 37~s。原因是独立方案的 42~s 会造成 8~s 逾期，第一优先级已经改变。机制因此具有明确的条件性：既不会因额外 2~s 工作一概拒绝，也不会因 R1 能提前就一概接受。

\subsection{独立执行与共享资源协调}\label{sec:3.7}

在 Gazebo 历史实验中，八个实体分别运行在独立进程中，各自保留任务状态和可选方案。实体通过消息与队友、任务站及入口服务站协调，由相应服务站统一管理共享资源。这样可以在同一轮试验中观察两件事：实体何时根据自己的任务作出选择，以及这项选择何时得到资源管理方许可。

这一回合记录了 17708 次跨进程请求。该数字反映仿真中的接口调用量。共享资源由服务站协调，参考位姿由仿真环境提供。

同一申请使用唯一编号，重复投递保留其状态，避免同一次请求产生两次占用。消息记录采用统一的仿真时基，可据此核对发送、到达和重试的顺序。第~\ref{sec:4.9} 节在 B3 到运输车的观察链路上注入随机丢包和断链。无线射频传播与实物网络需另测。

Webots 整场由裁判统一核对通行条件，各实体根据各自的传感与运动状态执行。集中核对有助于避免两个入口同时判空，独立执行则保留了单车对局部状态的响应。

\subsection{协同行为与实验小结}\label{sec:3.8}

\begingroup\small
\setlength{\tableavailable}{\dimexpr\linewidth-6\tabcolsep\relax}
\begin{longtable}{@{}>{\raggedright\arraybackslash}p{0.21000\tableavailable}>{\raggedright\arraybackslash}p{0.31000\tableavailable}>{\raggedright\arraybackslash}p{0.22000\tableavailable}>{\raggedright\arraybackslash}p{0.26000\tableavailable}@{}}
\caption{多体协同行为与关键技术}\label{tab:19}\\
\toprule
\textbf{协同行为} & \textbf{关键技术} & \textbf{仿生参考} & \textbf{实验观察} \\
\midrule
\endfirsthead
\multicolumn{4}{@{}l}{\small\color{muted}表~\thetable\quad （续）}\\
\toprule
\textbf{协同行为} & \textbf{关键技术} & \textbf{仿生参考} & \textbf{实验观察} \\
\midrule
\endhead
\midrule
\multicolumn{4}{r@{}}{\footnotesize\color{muted}续下页}\\
\endfoot
\bottomrule
\endlastfoot
共享通道互斥通行 & 车体包络判空（包络半径 0.28~m）、串行放行、通道外对准 & 分离规则的局部间距思想~\cite{ref7} & 三组整机 0 重叠，按实际清空放行少等 24.112~s \\
信息时效与放行 & 报告带采样时刻，6~s 有效期，拒收过期帧并主动补采 & 少数知情个体引导群体~\cite{ref9} & 补采使按期数由 13/14 恢复到 14/14 \\
岗位协商与支援 & 提议、确认、前往、待命、执行、释放六态岗位，有限候选筛选 & 响应阈值分工~\cite{ref8} & 按需询问减少报价次数，各组按期率相同 \\
支援的机会成本 & 字典序效用向量：按期数、逾期、活跃工作、完成时刻 & 军蚁活桥的成本收益折中~\cite{ref6} & 12 组条件中 2 组选择改变 \\
资源预约与提前释放 & 名义时间窗预约、实际清空后重算释放、联合准入 & — & 提前释放增加可预约余量，交付收益还受目的地占用影响 \\
\end{longtable}
\endgroup
